\subsection{Задание}

Исследуйте пример из дополнительных материалов к Лекции 6: \texttt{UFSACO.ipynb}, данные \texttt{samples.pkl}, а также методические указания (\textit{Лекция 6. Комментарии к коду (FS).docx}). 

Данный кейс демонстрирует отбор значимых признаков при моделировании содержания меди в продуктах обработки ПАО «Норильский никель» (процесс пенной флотации). Вам необходимо решить аналогичную задачу отбора признаков для собственных данных, используя предложенный программный комплекс.

\subsubsection*{Порядок выполнения работы:}

\begin{enumerate}
    \item Выберите набор данных для задачи \textit{регрессии} на одном из агрегаторов задач машинного обучения (например, \href{https://www.kaggle.com/datasets}{Kaggle}, \href{https://archive.ics.uci.edu}{UCI Machine Learning Repository} и др.). 
    
    Определите вещественный целевой признак. Приведите в отчете краткое описание предметной области датасета, его структуры и признакового состава, опираясь на информацию с портала и результаты разведочного анализа данных.

    \item Используя алгоритм градиентного бустинга \textit{(в режиме «чёрного ящика»)}, определите подмножество наиболее значимых признаков для предсказания целевой переменной.

    \item С помощью алгоритма ненаправленного муравьиного отбора признаков --- UFSACO --- выделите значимые признаки \textit{без учета целевой переменной}. 
    
    \item Убедитесь, что пересечение двух множеств признаков, полученных методами бустинга и ACO, не является пустым. Мощность пересечения должна составлять \textbf{не менее 4 элементов}. Если данное требование не выполняется, поварьируйте гиперпараметры алгоритма муравьиной колонии.
\end{enumerate}

\paragraph{Важно!} Допускается использование только той версии алгоритма ACO, которая программно представлена в файле \texttt{UFSACO.ipynb}.

\paragraph{Ограничения и требования к датасету:}
\begin{itemize}
    \item Количество признаков: $\ge 20$.
    \item Количество объектов: $\ge 100$.
    \item Дата публикации набора данных: не ранее 2010 года.
    \item Количество отобранных признаков каждым из методов: не менее 7 и не более половины от общего числа признаков в исходном датасете.
    \item Во избежание дублирования выбранный датасет необходимо предварительно согласовать.
\end{itemize}